% Generated by experiments/export_paper_tables.py -- do not edit by hand.
% Requires \usepackage{booktabs}.  Wide tables use table* (IEEEtran two-column).

% ---- Table 1: Main accuracy (compact)
\begin{table*}[t]
\centering
\caption{Test accuracy, mean $\pm$ std over seeds (ddof=1): RMSE on the three CALPHAD splits and raw-unit MAPE on the temp.-band split. Bold = best mean in the random and temp.-band splits (no ranking in the comp.-group split).}
\label{tab:c_accuracy}
\scriptsize
\setlength{\tabcolsep}{3pt}
\begin{tabular}{lrrrr}
\toprule
 & \multicolumn{3}{c}{RMSE (J/mol)} & \multicolumn{1}{c}{MAPE (\%)} \\
\cmidrule(lr){2-4}\cmidrule(lr){5-5}
Method & Random & Temp.-band & Comp.-group & Temp.-band \\
\midrule
SymGWR-T & $2.25 \pm 0.15$ & 43.0 $\pm$ 0 & $124 \pm 98$ & 0.0369 $\pm$ 0 \\
SymGWR-R & $1.16 \pm 0.10$ & 2.24 $\pm$ 0 & $60.0 \pm 64.3$ & $1.54\times10^{-3}$ $\pm$ 0 \\
\midrule
SymGWR-R-CS (k=4) & $1.01 \pm 0.49$ & $3.65 \pm 0.06$ & $59.6 \pm 64.4$ & $(1.90 \pm 0.01)\times10^{-3}$ \\
Cubic along T / arm$^{a}$ & \textbf{\boldmath $0.800 \pm 0.448$} & \textbf{\boldmath 0.553 $\pm$ 0} & $3.34 \pm 2.36$ & \textbf{\boldmath $2.65\times10^{-4}$ $\pm$ 0} \\
\midrule
Polynomial & $18.1 \pm 0.9$ & 17.5 $\pm$ 0 & $53.3 \pm 17.9$ & 0.0120 $\pm$ 0 \\
GP$^{b}$ & $8.40 \pm 0.45$ & $8.42 \pm 0.10$ & $84.0 \pm 106.8$ & $(5.84 \pm 0.07)\times10^{-3}$ \\
MLP & $129 \pm 17$ & $246 \pm 29$ & $324 \pm 285$ & $0.212 \pm 0.024$ \\
PINN$^{c}$ & $223 \pm 134$ & $285 \pm 169$ & $309 \pm 218$ & $0.259 \pm 0.170$ \\
XGBoost & $154 \pm 8$ & $9079 \pm 2$ & $928 \pm 423$ & 8.66 $\pm$ $2.7\times10^{-3}$ \\
k-NN & $578 \pm 25$ & $4675 \pm 3$ & $635 \pm 270$ & 4.78 $\pm$ $3.8\times10^{-3}$ \\
LOESS$^{d}$ & $425 \pm 11$ & $608 \pm 7$ & $1087 \pm 1746$ & $0.566 \pm 0.007$ \\
\bottomrule
\end{tabular}
\par\smallskip
\begin{minipage}{\linewidth}\scriptsize\raggedright
"$\pm$ 0": std below 1e-08 of the mean (numerical noise; fixed test set and deterministic fit), not exactly zero.\par
SymGWR-T = degree-5 polynomial + nearest-reference Taylor surrogate; SymGWR-R = residual-anchored; SymGWR-R-CS = residual-anchored with a compactly supported blend of the k=4 nearest reference residuals; cubic = interpolant along T (random, temp.-band) or along the composition arm (comp.-group).\par
Effective independent partitions: Random 10 (10 random 60/20/20 resamples; test sets overlap, pairs not strictly independent); Temp.-band 1 (test set fixed (1,568 rows, same for every seed); only train/val partition varies); Comp.-group 5 (5 folds x 2 shuffles: test fold = seed \% 5, so 5 distinct test sets, each seen twice).\par
(a) Comp.-group: cubic along the composition arm is defined for only 81\% of the test rows; its metrics are on those rows only and are not directly comparable with the other rows.\par
(b) GP protocol: alpha tuned on a 1000-row train subsample (val MSE), then ONE refit on train (capped 2500) + val; 25 min timeout\par
(c) temperature column moved to index 0 before train\_pinn: the library physics loss penalises dG/d(feature 0), and feature 0 in the data is Mass\_fraction\_Nb\par
(d) train rows capped at 2500 (random subsample, seed-dependent); val rows added on refit\par
Comp.-group: no method is marked best. Paired t-test on log RMSE and log MAPE over the 5 test folds (10 seeds = 5 folds $\times$ 2 shuffles, averaged within a fold; 4 df); meaningful = same sign in all 5 folds and p $<$ 0.05 on both metrics, uncorrected; cubic row not tested (different test rows). Leading group (Polynomial, SymGWR-R-CS, SymGWR-R, GP, SymGWR-T): differences within it are within fold variation, except GP and SymGWR-T are worse than SymGWR-R-CS, SymGWR-R. The other methods shown (LOESS, MLP, PINN, XGBoost, k-NN) are worse than every member of the leading group.\par
\end{minipage}
\end{table*}

% ---- Table 2: Derivative fidelity (compact)
\begin{table}[t]
\centering
\caption{Analytic gradient of the degree-5 SymGWR polynomial vs. central finite differences (FD) of the raw CALPHAD table on all grid points. Normalized error = MAE / typical \textbar{}derivative\textbar{}; r = Pearson correlation; mean $\pm$ std over 10 seeds (ddof=1).}
\label{tab:c_derivfid}
\scriptsize
\setlength{\tabcolsep}{3pt}
\begin{tabular}{lrrr}
\toprule
First derivative & n & Norm. err. (\%) & r \\
\midrule
dG/dT [J/(mol K)] & 7742 & $0.0531 \pm 0.0017$ & 0.9999 \\
dG/dNb & 1944 & $1.32 \pm 0.01$ & 0.9986 \\
dG/dAl & 648 & $1.55 \pm 0.04$ & 0.9997 \\
dG/dTi & 1539 & $1.34 \pm 0.01$ & 0.9996 \\
dG/dMo & 1539 & $3.81 \pm 0.01$ & 0.9978 \\
dG/dFe & 648 & $9.89 \pm 0.44$ & 0.9907 \\
dG/dCr & 729 & $2.77 \pm 0.10$ & 0.9868 \\
dG/dNi & — & — & — \\
\midrule
Second derivative & n & vs narrow FD (\%) & vs best ref. (\%) \\
\midrule
d$^2$G/dT$^2$ & 7742 & $10.9 \pm 0.1$ & $2.59 \pm 0.10$ \\
\bottomrule
\end{tabular}
\par\smallskip
\begin{minipage}{\linewidth}\scriptsize\raggedright
Fit design rank 95 of 1286 polynomial terms; composition-axis accuracy holds for own-axis derivatives on the sweep lines only (identified there; other planes not identified, null share 0.71--0.97). Ni is a constant column: no FD, not identifiable (null share 1.00).\par
d$^2$G/dT$^2$: the narrow FD (one grid step) is noise-limited (noise share 83\% of its squared error, from the rounding of the table), so the error vs the best reference (T$\pm$5 wide-stencil FD or smoothing spline) is the one to read.\par
Pure d$^2$G/dc$^2$ for Fe and Cr: model error 32--42\% (other axes 5--21\%). Mixed d$^2$G/dT dc: noise-limited in the narrow FD for 4 of 6 axes (not tabulated). Cross-element mixed terms (30) are not identified (null share 0.91--1.00) and have no FD reference.\par
\end{minipage}
\end{table}

% ---- Table 3: Continuity at nearest-reference switches (compact)
\begin{table}[t]
\centering
\caption{Jump of the surrogate when the nearest reference changes, on 5,000 short segments per split and seed (10 seeds); median and p95 averaged over seeds. Value jump in J/mol; gradient jump relative to the median model gradient norm.}
\label{tab:c_continuity}
\scriptsize
\setlength{\tabcolsep}{3pt}
\begin{tabular}{lrrr}
\toprule
Variant & Value, median & Value, p95 & Gradient, median \\
\midrule
\multicolumn{4}{l}{\textit{Random split}} \\
\midrule
SymGWR-T & 0.480 & 4.92 & 0.0248 \\
SymGWR-R & 0.477 & 3.96 & $6.17\times10^{-17}$ \\
SymGWR-R-CS (k=4) & $2.53\times10^{-12}$ & $1.01\times10^{-11}$ & $3.75\times10^{-7}$ \\
\midrule
\multicolumn{4}{l}{\textit{Temp.-band split}} \\
\midrule
SymGWR-T & 58.4 & 422 & 0.123 \\
SymGWR-R & 3.59 & 24.5 & $1.42\times10^{-17}$ \\
SymGWR-R-CS (k=4) & $1.74\times10^{-14}$ & $2.17\times10^{-12}$ & $2.56\times10^{-6}$ \\
\midrule
\multicolumn{4}{l}{\textit{Comp.-group split}} \\
\midrule
SymGWR-T & 9.32 & 190 & 0.0134 \\
SymGWR-R & 0.366 & 19.5 & $7.83\times10^{-17}$ \\
SymGWR-R-CS (k=4) & $2.08\times10^{-12}$ & $1.21\times10^{-11}$ & $1.62\times10^{-16}$ \\
\bottomrule
\end{tabular}
\par\smallskip
\begin{minipage}{\linewidth}\scriptsize\raggedright
NN switch = the nearest reference changes. SymGWR-T extrapolates linearly from the nearest reference, so value and gradient both jump. SymGWR-R adds a constant offset to the polynomial, so its gradient is continuous (jumps are zero up to rounding) and only the value jumps. SymGWR-R-CS (k=4) blends the residuals of the k nearest references with compact support; its jumps at NN switches are small in value and gradient.\par
\end{minipage}
\end{table}

% ---- Table 4: Newton test (compact)
\begin{table}[t]
\centering
\caption{Implicit cooling of one lumped body per composition (98 compositions, 150 steps) with each surrogate's analytic Jacobian dH/dT = $-$T d$^2$G/dT$^2$; mean over 10 seeds $\times$ 98 compositions (980 runs) per surrogate and protocol. Converged = every step reaches \textbar{}R\textbar{} $<$ tol within 20 iterations; escaped = the accepted temperature left the data window.}
\label{tab:c_newton}
\scriptsize
\setlength{\tabcolsep}{3pt}
\begin{tabular}{llrrr}
\toprule
Surrogate & Protocol & Converged (\%) & Mean iters & Escaped (\%) \\
\midrule
Reference (truth) & --- & 100 & 2.89 & 0 \\
Polynomial & Random & 100 & 2.69 & 0 \\
Polynomial & Temp.-band & 100 & 2.69 & 0 \\
SymGWR-R & Random & 91.1 & 2.74 & 0 \\
SymGWR-R & Temp.-band & 93.9 & 2.75 & 0 \\
SymGWR-T (Newton) & Random & 0 & 20.0 & 100 \\
SymGWR-T (Newton) & Temp.-band & 0 & 20.0 & 100 \\
Spline (train rows) & Random & 75.2 & 3.65 & 17.8 \\
Spline (train rows) & Temp.-band & 56.1 & 4.11 & 27.6 \\
Spline (full raw table) & --- & 53.1 & 4.19 & 30.6 \\
SymGWR-R-IDW (k=4) & Random & 9.08 & 7.74 & 53.2 \\
SymGWR-R-IDW (k=4) & Temp.-band & 4.59 & 8.83 & 66.7 \\
SymGWR-R-CS (k=4) & Random & 8.78 & 7.84 & 56.2 \\
SymGWR-R-CS (k=4) & Temp.-band & 0.408 & 9.90 & 68.7 \\
\bottomrule
\end{tabular}
\par\smallskip
\begin{minipage}{\linewidth}\scriptsize\raggedright
Truth = degree-4 smoothing spline of G (the raw table is rounded, so second derivatives of an interpolating spline are quantisation noise). Reference and full-raw-spline rows run once per composition (98 runs, no protocol). Spline (train rows) = cubic spline along T through the fit rows of the composition. Secant and Brent solvers for SymGWR-T are not shown.\par
\end{minipage}
\end{table}

% ---- Table 5: Query time vs. number of reference points (compact)
\begin{table}[t]
\centering
\caption{Mean per-query latency ($\mu$s) over 10,000 timed queries, averaged over 3 repeats. Last row: log-log slope of query time against n with 95\% CI (slope 0 = constant time).}
\label{tab:c_scaling}
\scriptsize
\setlength{\tabcolsep}{3pt}
\begin{tabular}{lrrr}
\toprule
n & SymGWR-T & SymGWR-R & Cubic spline \\
\midrule
1,000 & 19.3 & 55.3 & 23.0 \\
6,350 (real) & 20.3 & 56.5 & 23.5 \\
100,000 & 93.2 & 135 & 28.6 \\
1,000,000 & 268 & 312 & 33.6 \\
\midrule
Slope [95\% CI] & 0.41 [0.34, 0.48] & 0.26 [0.21, 0.31] & 0.06 [0.05, 0.07] \\
\bottomrule
\end{tabular}
\par\smallskip
\begin{minipage}{\linewidth}\scriptsize\raggedright
Sizes above 6,350 are tiled copies of the real reference points with 1\% jitter: they test index (KD-tree) behaviour, not accuracy. n = 6,350 is the real fit. The slope rests on 4 distinct sizes $\times$ 3 repeats. Timing: process pinned to core 10, OMP\_NUM\_THREADS=1, GC disabled while timing; cubic spline = scipy.\par
Speedup vs pycalphad calculate, precompiled phase records: mean latency 0.8735 ms over 100 queries, i.e. 43$\times$ SymGWR-T and 15$\times$ SymGWR-R at n = 6,350. The pycalphad database is a steel database, not the system SymGWR is fitted to, and pycalphad was timed in a separate run: orders of magnitude only.\par
\end{minipage}
\end{table}
